\section{Segmentation by Thresholding}

Using \texttt{numpy} one can easily make the following function which takes an
RGB image in either grayscale or color (1 or 3 channels) and a threshold value
to produce a binary segmentation mask.

\begin{lstlisting}
def threshold_image(
    image: NDArray[Any | np.uint8],
    threshold: int,
) -> NDArray[np.uint8]:
    assert threshold >= 0 and threshold <= 255, "threshold must be in [0,255] range"
    if image.ndim == 3 and image.shape[2] == 3:
        gray = rgb2gray(image)*255 # rescale from [0,1] to [0,255]
    else:
        gray = image.astype(np.uint8)

    binary_mask = np.where(gray >= threshold, 1, 0).astype(np.uint8)
    
    return binary_mask
\end{lstlisting}
It is worth noting, that by converting to grayscale with \texttt{skimage.color.rgb2gray},
we are using a weighted average which heavily favors green.
This weighting is supposed to reflect the color sensitivities of the human eye, which is most sensitive to green light.
For this particular segmentation task, it proves beneficial, but when the objective is not
a grayscaling which corresponds more accurately to our visual perception,
but rather background-foreground seperation,
it may be worth considering other methods of grayscaling, such as a simple average of the RGB channels.

By plotting the binary mask as an image with e.g. \texttt{matplotlib.pyplot.imshow} as grayscale,
the binary mask is automatically scaled such that 1 is a white pixel and 0 is a black one. The resulting image can be seen below:
\begin{figure}[H]
	\centering
	\includegraphics[width=0.9\textwidth]{images/thresholding.png}
	\caption{Original, grayscale, and binary segmentation of the same \texttt{pillsect.png} image sidy-by-side.}
	\label{fig:thresh}
\end{figure}

For this image, we see that we have somewhat succesfully seperated the foreground and the background.
The green object, when converted to grayscale, which is done by averaging the three RGB color channels.
We can see that the grayscale image has very little distinction from the original, except for the green object.
This green object becomes rather gray in the grayscale image, which in turn results in our thresholding function
falling somewhat short for seperating the object from the background perfectly.
However the other objects are very well seperated.

Another noticable result, is the salt-like noise. This is due to the texture of the background, 
having a decent splatter splatter of bright and light-gray to white pixels.
It would be interesting to see the result of applying some filtering kernel like e.g. Gaussian, mean, or median prior 
to thresholding. Depending on the task at hand, this may or may not be necessary.

Overall a very decent segmentation, albeit with quite a bit of salt-like noise.

\subsection{Intensity Histogram}

We now plot the (pixel) intensity histogram of the grayscale image.

\begin{figure}[H]
	\centering
	\includegraphics[width=0.65\textwidth]{images/histogram.png}
	\caption{Pixel intensity histogram of the grayscale image.}
	\label{fig:histogram}
\end{figure}

The plot includes a dotted red vertical line indicating where we set the threshold.
From the histogram, it is clear that we successfully segmented the two major clusters of pixel intensities:
the cluster between $\approx 20$ and $\approx 90$, and the cluster of white pixels around 255.
We do, however, observe that some pixels are being cut off which arguably belong to the low-intensity cluster.
This "tail" of higher intensities is most likely part of the background,
and assuming this is the case, setting the threshold somewhat higher 
-- to, say, 125 -- could have resulted in less salt-like noise.
From the histogram alone, however, we cannot distinguish which frequencies in the 
$[100, 150]$ interval correspond to salt-like noise and which belong to low-intensity green object pixels,
so the current threshold of 100 is not necessarily incorrect.
